% =============================================================================
% Title Page
% =============================================================================

\begin{titlepage}
\raggedright

% Title
{\fontsize{36pt}{40pt}\selectfont\bfseries\color{OECDblue}
Estimating Climate Change Risk Exposure: a probabilistic approach\par}

\vspace{2cm}

% Authors
\raggedleft
\setlength{\parskip}{0.3em}

\textbf{Camilo Saldarriaga}\\
ESCP Business School\\
\href{mailto:juan_camilo.saldarriaga@edu.escp.eu}{juan_camilo.saldarriaga@edu.escp.eu}

\vspace{1.5em}

\textbf{Jesse Grabowski}\\
Paris-1 Panth\'{e}on-Sorbonne\\
\href{mailto:jessegrabowski@gmail.com}{jessegrabowski@gmail.com}

\vspace{1.5em}

\textbf{Jan Rielaender}\\
OECD Development Center\\
\href{mailto:Jan.RIELAENDER@oecd.org}{Jan.RIELAENDER@oecd.org}

\vspace{2cm}

% Abstract
\raggedright
\textbf{Abstract}

\justifying
This paper develops a methodology to quantify countries' exposure to climate change--related disasters, focusing on hydrometeorological events such as floods and storms. We introduce the Probabilistic Geospatial Risk (PGR) model, a fully Bayesian, multi-level system integrating time series, frequency, damage, and spatial sub-models. The framework captures uncertainty across the entire distribution of disaster frequency and damages, enabling robust estimation of both mean and tail risks. Using data from the EM-DAT disaster database, the World Bank, GPCC, and NOAA, we estimate disaster probabilities and damage distributions as functions of climate and development indicators. We further apply a Hilbert Space Gaussian Process (HSGP) to model fine-grained geospatial patterns of disaster occurrence within countries, complemented by state-space projections of climate and economic covariates. The model is demonstrated through two case studies---Costa Rica and the Lao People's Democratic Republic. The resulting geolocalized damage curves and return-year estimates provide a transparent, reproducible tool for assessing future exposure to climate-related risks. This probabilistic approach advances the literature by integrating uncertainty quantification, geospatial modeling, and economic impact estimation into a coherent framework suitable for policy analysis and climate adaptation planning.

\vspace{1.5em}

\textbf{JEL Classification:} Q54, Q51, C11, C150, C53, Q56, and D81

\vspace{0.5em}

\textbf{Keywords:} Climate change risk, Bayesian modeling, hydrometeorological disasters, geo-spatial modeling, probabilistic risk assessment, damage curves, uncertainty quantification.

\end{titlepage}
